\problem{MEC-0036}{36}
(IJSO รอบที่ 1 2553) อนุภาคหนึ่งถูกโยนขึ้นในแนวดิ่งภายใต้สนามโน้มถ่วง $g$ กราฟความเร็ว--เวลาเป็นดังรูป (ทิศขึ้นแทนด้วยเครื่องหมายบวก) ตอนขาลงเมื่ออนุภาคมีอัตราเร็วครึ่งหนึ่งของอัตราเร็วต้น เวลาผ่านไปเท่าใดนับจากเริ่มโยน
\begin{center}\begin{tikzpicture}[>=Stealth]\draw[->] (0,-1.1)--(0,2.4) node[above]{ความเร็ว};\draw[->] (-.3,0)--(5,0) node[right]{เวลา};\draw[thick] (0,1.6)--(4.4,-1.1);\node[below left] at (0,0){$0$};\node[below] at (2.61,0){$T$};\end{tikzpicture}\end{center}
\par\smallskip\par\smallskip\begin{tabular}{@{}p{.48\linewidth}p{.48\linewidth}@{}}1. $\dfrac{T}{2}$ & 2. $T\left(1+\dfrac{1}{\sqrt{2}}\right)$ \\ 3. $\dfrac{3T}{2}$ & 4. $T\left(2-\dfrac{1}{\sqrt{2}}\right)$\end{tabular}
